\documentclass[11pt,letterpaper]{article}
\usepackage[T1]{fontenc}
\usepackage{lmodern}
\usepackage[margin=1in]{geometry}
\usepackage{amsmath,amssymb,amsthm,mathtools}
\usepackage{booktabs,tabularx,longtable,array}
\usepackage{microtype}
\usepackage{xcolor,xurl,hyperref}
\usepackage{enumitem,fancyhdr}
\usepackage{listings}
\definecolor{ink}{HTML}{18334B}
\definecolor{linkblue}{HTML}{185B8D}
\hypersetup{colorlinks=true,linkcolor=linkblue,urlcolor=linkblue,citecolor=linkblue,
  pdftitle={K215: Checked components for shorter correctness proofs},
  pdfauthor={OpenAI, prepared for Vladimir Reshetnikov}}
\pagestyle{fancy}
\fancyhf{}
\fancyhead[L]{\small\textcolor{ink}{K215: contextual proof simplification}}
\fancyhead[R]{\small 22 October 2026}
\fancyfoot[C]{\thepage}
\setlength{\headheight}{14pt}
\setlength{\parskip}{5pt}
\setlength{\parindent}{0pt}
\setlist{itemsep=3pt,topsep=5pt}
\lstset{basicstyle=\ttfamily\small,columns=fullflexible,breaklines=true,
  frame=single,rulecolor=\color{black!15},backgroundcolor=\color{black!2},
  xleftmargin=4pt,xrightmargin=4pt,showstringspaces=false}
\newtheorem{theorem}{Theorem}[section]
\newtheorem{lemma}[theorem]{Lemma}
\newtheorem{corollary}[theorem]{Corollary}
\newtheorem{proposition}[theorem]{Proposition}
\theoremstyle{definition}
\newtheorem{definition}[theorem]{Definition}
\newtheorem{remark}[theorem]{Remark}
\newcommand{\code}[1]{\texttt{\detokenize{#1}}}
\newcommand{\Con}{\operatorname{Con}}
\newcommand{\Halts}{\operatorname{Halts}}
\newcommand{\rev}{\operatorname{rev}}
\newcommand{\flatten}{\operatorname{flat}}
\newcommand{\conc}{\operatorname{conc}}
\newcommand{\frameop}{\operatorname{frame}}
\newcommand{\sstep}{\operatorname{sstep}}
\newcommand{\Reach}{\operatorname{Reach}}
\newcommand{\ReachOne}{\operatorname{Reach1}}
\newcommand{\Some}{\operatorname{Some}}
\newcommand{\None}{\operatorname{None}}
\newcommand{\runlocal}{\operatorname{runLocal}}
\newcommand{\RegEq}{\mathrel{\simeq}}
\newcommand{\baseurl}{https://github.com/VladimirReshetnikov/ProveIt/blob/cbf94cb552a3b49509e0eee80e7455d4eb5974c6/}

\begin{document}
\hypersetup{pageanchor=false}
\begin{titlepage}
\vspace*{0.2in}
{\Large\color{ink} ProveIt research note\par}
\vspace{0.25in}
{\huge\bfseries K215: Checked Components for\\[5pt] Shorter Correctness Proofs\par}
\vspace{0.18in}
{\Large Guarded control frames and extensional\\source rewriting in Lean and Rocq\par}
\vspace{0.3in}
{\large Prepared for Vladimir Reshetnikov\par}
22 October 2026
\vspace{0.3in}
\begin{abstract}
The accepted 215-state binary Turing machine already has substantial proof
simplifications. The next useful step is to prove reusable statements about
unchanged context, then instantiate them at repeated occurrences. This note
turns two existing research proposals into checked components.

In Lean 4.32.0, we prove a guarded append-left frame rule for the maintained
symbolic executor, check all 35 proposed control templates, and derive complete
whole-tape and source-action contracts for every template. An independent audit
binds those templates to 1,422 control occurrences of the exact current source.
The proposed mixed proof has 4,789 expensive replay obligations instead of
6,176, while retaining all 6,176 source-outcome bindings. This is an obligation
count, not a measured endpoint speedup or a completed replacement endpoint.
An independent Rocq port checks the same 35 local replays and their framed
reachability; its Home and source-action adapters remain integration work.

For the Rocq source proof, we introduce pointwise register transport, avoiding
functional extensionality, and prove that reversing exactly shared nodes 188
and 190 of \code{dagPerm} preserves all structured exits and iterated halting.
These components are checked with Coq 8.18.0 and recursive \code{coqchk}; the
production Rocq 9.1.0 endpoint is a separate integration target.

The report gives mathematical proofs, exact integration obligations, failure
examples for unsafe simplifications, provenance, and further research
questions. It does not change the literal machine, claim an unconditional
consistency result, or claim a fresh build of either full endpoint.
\end{abstract}
\vfill
\textbf{Repository snapshot}\par
\small\nolinkurl{cbf94cb552a3b49509e0eee80e7455d4eb5974c6}\par
\normalsize\textbf{Primary deliverable:} checked reusable proof components,
self-contained replay sources, and a staged integration proposal.
\end{titlepage}
\hypersetup{pageanchor=true}

\tableofcontents
\newpage

\section{Recommendation and scope}

The recommended approach is \emph{contextual proof reuse}: isolate a local
transformation, state exactly what surrounding data it cannot inspect or alter,
and prove the transformation once for every admissible surrounding context.
For K215, this is useful in two different layers:

\begin{enumerate}
\item In the binary runtime proof, a short control routine changes a suffix of
the encoded program counter. Its earlier prefix is preserved until the
existing Home routine scans it and returns to dispatch. One quantified
contract can therefore replace many repeated symbolic executions.
\item In the source proof, a node whose body increments two registers can be
replaced by the opposite order. The surrounding program sees the same register
values and the same structured exit. A generic contextual transport theorem
therefore propagates two local proofs through all their occurrences.
\end{enumerate}

These ideas were already identified in two repository proposals:
\code{round13-left-frame} \cite{leftproposal} and
\code{k215-simplification-3} \cite{cfgproposal}.
Their discovery is not claimed here.
The contribution of this delivery is to supply missing universal proofs and
checked concrete instances, including an extensional-register interface that
preserves Rocq's closed assumption boundary.

The most valuable near-term integration is the 35-template runtime family.
Its generic theorem, exact literal executions, right-tape normalization, and
source control actions have now been checked. The remaining work is finite
source-occurrence assembly, semantic coverage, and endpoint rechecking.
The two-node source rewrite is a second, complementary component. It still
needs a certified source-to-CFG connection before it can replace the current
source correspondence certificate.

The final theorem remains
\[
 \Halts(K215,\mathrm{blank})
 \quad\Longleftrightarrow\quad
 \text{ZFC has a finite derivation of falsity}.
\]
Its negation yields the intended syntactic consistency equivalence. A proof
of this equivalence does not assert that K215 fails to halt unconditionally.
The repository distinguishes its external Lean/Rocq theorems from a separately
encoded derivation in PA or in ZFC's own proof syntax \cite{verification}.
This delivery preserves that distinction.

\section{What is already simplified}
\label{sec:baseline}

\subsection{Use current source, not a historical checkpoint count}

The review is pinned to commit
\nolinkurl{cbf94cb552a3b49509e0eee80e7455d4eb5974c6}.
The phase-15 checkpoint records an endpoint with 176 project modules and
2,052,142 selected source bytes \cite{phase15}. An independent recursive import
walk of the maintained \code{KEndpoint215} at the review pin instead finds
\textbf{174 project modules, 1,752,650 raw source bytes, and 37,648 lines}, with
\code{Init} as the only external import boundary. The maintained reproduction
driver's inventory routine gives the same census.

This is a \emph{static current-source measurement}. It is not a new compilation
of those 174 modules, and historical phase-15 receipts must not be represented
as fresh checks of all current bytes. The source census and its exact hashes
are included in the archive. Later direct-seed/interface and source-calculus
changes explain why the maintained closure can differ from the checkpoint
summary. No claim that the current endpoint is invalid follows from this
documentation lag.

\begin{table}[ht]
\centering\small
\begin{tabularx}{\textwidth}{@{}lrrX@{}}
\toprule
Scope & Modules & Source bytes & Interpretation\\
\midrule
Phase-15 endpoint & 176 & 2,052,142 & Historical checkpoint and its receipts\\
Current import closure & 174 & 1,752,650 & Fresh static census at the review pin\\
New component checks & \multicolumn{2}{l}{See the included build records} & Fresh focused compilation; no full endpoint replay\\
\bottomrule
\end{tabularx}
\caption{Different evidence scopes must not be conflated.}
\end{table}

\subsection{Do not repeat completed work}

The repository already uses offset-bisimulation symmetry, normalized source
correspondence, shared dominated-implication arguments, a generic AX word
composition, shared enumerator contracts, smaller semantic import boundaries,
shared recursion and quotient syntax, and a compact dense-machine invariant
\cite{simplifications}. The source correspondence already has 1,426 normalized
checkpoint entries, rather than the older 2,240. The current source also
contains direct compilation of the enumerated calculus and the updated seed
route. Recommending these as entirely new changes would miss the maintained
baseline.

The final endpoint itself is already a short composition. Shortening its few
lines further would have little effect on the cost of checking the repeated
runtime and source certificates. We instead target statements that account
for repeated semantic work.

\subsection{Preserve the exact object and proof boundary}

The accepted object has 215 working states, two symbols, a two-way infinite
tape, blank input, initialization included, and a separate halt. The dense
table has 430 transition rows \cite{candidate,verification}.
The following identities are retained:

\begin{description}[style=nextline]
\item[Dense table SHA-256]
\small\nolinkurl{fdbcef0825886035c3091c8537a69535bc96f9e265f0ae45863543d40ae70090}
\item[Native table SHA-256]
\small\nolinkurl{9683d6ed9628e6ce407c315b08f5da3d063da94e92706f275ca0f33e8272e4b2}
\item[Source SHA-256]
\small\nolinkurl{17af207195179b2508bf91154cd00739374f8862207ca21b9a3cb9098aa6f5ba}
\end{description}

\normalsize The fresh focused Lean proofs use ordinary kernel-checked
\code{decide}, not \code{native_decide}, and contain no accepted \code{sorry}
or added axioms. Their reports stay within the existing Lean allowance
\code{propext}, \code{Classical.choice}, and \code{Quot.sound}.
The fresh Rocq source component reports are closed under the global context.
Conditional hypotheses such as the four existing register-routine contracts
are explicit theorem parameters, not postulated axioms.

\section{Where the runtime repetition comes from}

The current source DAG has 202 distinct nodes and expands to 4,751 leaves.
The independent audit follows the maintained 234-role decoder and compares
all 430 formal table rows with the native table after the explicit label
translation. The leaf and guarded-outcome counts are:

\begin{table}[ht]
\centering
\begin{tabular}{@{}lrr@{}}
\toprule
Primitive kind & Source leaves & Guarded outcomes\\
\midrule
Increment & 1,900 & 1,900\\
Decrement/test & 1,425 & 2,850\\
Break & 98 & 98\\
Continue & 1,327 & 1,327\\
Halt & 1 & 1\\
\midrule
Total & 4,751 & 6,176\\
\bottomrule
\end{tabular}
\caption{A decrement contributes separate zero and nonzero outcomes.}
\end{table}

The framed family covers 1,326 continue occurrences and 96 break occurrences:
1,422 control cases in total. It uses exactly 35 distinct local bodies. The
three control fallbacks are \code{001/CB1}, \code{1000111/CC3}, and
\code{11011/CB1}. The 35 bodies contain 444 symbolic operations altogether,
with a maximum of 39 in one body. All 35 bodies are used by the audited
occurrence map.

The proposed expensive replay census is therefore
\[
 35+3+1900+2850+1=4789,
\]
compared with 6,176 literal outcome replays. The difference is
\[
 6176-4789=1387,
 \qquad \frac{1387}{6176}\approx22.4579\%.
\]
All 6,176 exact source-outcome bindings remain necessary. A template instance
still needs a route, suffix decomposition, prefix-scanner premise, action
identity, and gap arithmetic. This count does not measure total proof size,
object size, memory, or wall-clock time. It also does not justify removing any
outcome from source coverage.

Native state labels are not the maintained formal labels. In particular,
native 301, 302, and halt 254 become formal 102, 120, and 230. State 302 occurs
in the actual local traces, so this translation cannot be silently omitted.
The new Lean finite equations are evaluated against the maintained
\code{KTable215.tm215}; the independent audit also checks the complete row
translation.

\section{A guarded frame theorem for the actual executor}
\label{sec:frame}

\subsection{Symbolic configurations}

The maintained executor describes a symbolic configuration by
\[
 s=(q,o,L,R),
\]
where $q$ is a state, $o$ indicates which list carries the head, and $L,R$ are
lists of symbolic segments. Segments are a single Boolean cell, a constant
run of length $an+c$, or a register-file box. The interpretation
$\conc_{n,v}(s)$ is the complete canonical zipper configuration for gap
parameter $n$ and register list $v$.

For any segment list $F$, define
\[
 \frameop_F(q,o,L,R)=(q,o,L\mathbin{+\!+}F,R).
\]
This appends an untouched left context. It does not trim or normalize the
left list. Define the guard
\[
 \operatorname{safe}(q,o,L,R)
 \quad\Longleftrightarrow\quad \neg(o=\mathrm{true}\ \wedge\ L=[]).
\]

\begin{theorem}[Guarded append-left naturality]
\label{thm:naturality}
For every table $M$, branch condition $b$, segment list $F$, and symbolic
configuration $s$, if $s$ is safe, then
\[
 \sstep_{M,b}(\frameop_F(s))
 =\operatorname{map}\bigl((t,m)\mapsto(\frameop_F(t),m),\,
                         \sstep_{M,b}(s)\bigr).
\]
The equation preserves both successful outputs and failure, and preserves the
executor's progress Boolean.
\end{theorem}

\begin{proof}
Destructure the head side and the leading segment. On a right-side bit, a
machine step either prepends the written bit to $L$, leaves $L$ unchanged, or
fails. All three operations commute with appending $F$. A left-side bit
reorientation removes the first bit from a nonempty $L$; its tail becomes
$(\operatorname{tail}L)\mathbin{+\!+}F$. For a constant run, the whole-run
self-loop branch prepends or transfers the same run, and the two run-peeling
helpers commute with append by inspection of their natural-number cases.

The register-box branch leaves the left list unchanged or prepends one fixed
bit. The oracle result is independent of $F$. Unsupported symbolic shapes
remain unsupported. The only other branch is a left-facing empty list; it is
excluded by the guard. These are all branches of the maintained \code{sstep}.
The checked Lean theorem is \code{KLeftFrame.sstep_natural}.
\end{proof}

\begin{remark}[Why the guard is essential]
For $s=(q,\mathrm{true},[],R)$, the original executor inserts an implicit blank
when reorienting. After appending $F=[\operatorname{bit}(1)]$, it instead reads
that explicit one. Hence the unguarded equation is false, even for a table
with no machine transitions. The package contains a checked Lean
counterexample. The guard must hold before \emph{every} local step, not only
at the initial configuration.
\end{remark}

\subsection{Lift a finite body once, for all frames}

Define a bounded local replay by
\[
\begin{split}
 \runlocal_{M,b}(0,s)&=\Some(s),\\
 \runlocal_{M,b}(j+1,s)&=
 \begin{cases}
 \runlocal_{M,b}(j,t),&\text{if $s$ is safe and $\sstep_{M,b}(s)=\Some(t,m)$},\\
 \None,&\text{otherwise}.
 \end{cases}
\end{split}
\]
It is permitted to \emph{end} with an empty left list. It cannot execute the
next, unsafe boundary reorientation.

\begin{theorem}[Local replay under every appended frame]
\label{thm:runlocal}
Assume the existing register-routine hypotheses \code{BoxLemmas M}. If
$\runlocal_{M,b}(j,s)=\Some(t)$, then, for every $F,n,v$ satisfying the existing
branch-validity predicate,
\[
 \Reach_M\bigl(\conc_{n,v}(\frameop_F(s)),
              \conc_{n,v}(\frameop_F(t))\bigr).
\]
\end{theorem}

\begin{proof}
Induct on $j$. At zero the successful equation identifies the endpoints, so
reflexivity applies. At a successor, successful replay exposes a safe
configuration, a successful symbolic step, and a shorter replay. Apply
Theorem~\ref{thm:naturality}; apply the existing \code{sstep_sound} theorem to
that framed step; then compose its concrete reachability with the induction
hypothesis. No additional tape interpretation or trusted evaluator is used
\cite{runtime}.
\end{proof}

The conclusion is reflexive-transitive reachability, not positive
reachability. Some symbolic operations are reorientations or zero-length run
normalizations. Requiring positive progress at every symbolic step would be
incorrect. The existing final Home theorem provides the necessary positive
machine progress for the whole source action.

\subsection{Finish with the existing Home theorem}

All 35 local bodies end at $(2,\mathrm{true},[],R')$. Choose
\[
 F=\operatorname{map}(\operatorname{bit},\rev(P))\mathbin{+\!+}[\operatorname{bit}(1)].
\]
The caller supplies the existing scanner premise
\[
 \code{scanBound}\bigl(0,\operatorname{map}(\operatorname{bit},\rev(P))\bigr)
 =\Some(k).
\]
The already proved \code{scanBound_sound} gives the left-word scan across
$\rev(P)$. The final one-cell marker resets its counter. Applying the existing
\code{home_word} then returns to state 84 with left marker $[1]$ and right word
$P\mathbin{+\!+}\flatten_{n,v}(R')$ \cite{home}. Thus:

\begin{corollary}[Positive local-plus-Home contract]
\label{cor:home}
Assume \code{BoxLemmas M}, \code{KHome.Rows M}, and
\code{KHome.Boundary M}.
If $\runlocal_{M,b}(j,s)=\Some(2,\mathrm{true},[],R')$, then for every accepted
prefix $P$, gap $n$, and valid register list $v$,
\[
 \ReachOne_M\left(
 \conc_{n,v}(\frameop_F(s)),\;
 z_c(84,[1],P\mathbin{+\!+}\flatten_{n,v}(R'))\right).
\]
\end{corollary}

\begin{proof}
Compose Theorem~\ref{thm:runlocal} with the positive Home execution using
\code{Reach.trans_reach1}. The explicit empty local left list makes the Home
entry exactly the caller's frame; no new theorem about crossing an opaque
boundary is required.
\end{proof}

There is no condition that $P$ end in one. In the actual occurrence map, 721
accepted prefixes end in zero. Adding an ends-in-one premise would needlessly
lose sharing. There \emph{is} a blank far-left exterior premise inherited from
the zipper/Home boundary semantics. The generic algebraic frame theorem is
valid for arbitrary segment frames, but that does not turn Home into an
arbitrary-nonblank-exterior theorem.

\subsection{Independent Rocq check of the local frame rule}

The archive also ports the guarded rule to the maintained Rocq
\code{Runtime230Core} semantics. The effective Rocq table agrees with Lean on
all 430 actions after the documented native-label translation. The same
35 literal templates are checked independently with \code{vm_compute}:
444 symbolic steps in total, with a maximum local length of 39.

\code{Refactor.LeftFrame} proves the append-frame naturality theorem, its
essential-guard counterexample, and finite local replay soundness under the
four inherited box-routine hypotheses.

The template module \code{Refactor.LeftFrame215} proves
\code{all_checked} and a universally quantified \code{template_replay}:
every listed local replay remains valid for every symbolic left frame.
Composing it with the generic soundness theorem gives the corresponding
concrete reachability contract. The four generic and two concrete assumption
reports are closed under the global context.

Nine exact source files, including the guarded frame core, were freshly
compiled with Coq 8.18.0 and recursively checked with \code{coqchk}. The tenth
source, the generated template client, was then freshly compiled against
those exact dependency objects. Its receipt is separate: a second recursive
sweep including that client was not run. The included replay driver can
freshly compile and recursively check the complete ten-source component.

This Rocq result has a narrower scope than the Lean result below. It proves
the local replay and framed reachability, but does not yet port the Home scan
or the exact source-action adapter. It is therefore not a claim of 35 positive
complete Rocq control contracts or of a rebuilt Rocq endpoint.

\section{Connect runtime templates to exact source actions}
\label{sec:actions}

\subsection{Reject prefix escape explicitly}

The original path arithmetic uses least-significant-first operations internally.
For example, \code{stripRev b k} removes cells through the $k$th occurrence of
bit $b$, and \code{incRev} performs fixed-width increment. Their original
definitions remain unchanged.

The new \code{KControlFrame} module supplies two option-valued sufficient
checks. \code{stripChecked} fails if the specified suffix does not contain
enough matching bits. \code{incChecked} fails if a carry would leave the
suffix. Successful values certify exact arithmetic under every added prefix.

\begin{lemma}[Safe stripping and incrementing]
If $\code{stripChecked}(b,k,s)=\Some(t)$, then for every list $p$,
\[
 \code{stripRev}(b,k,s\mathbin{+\!+}p)=t\mathbin{+\!+}p.
\]
If $\code{incChecked}(s)=\Some(t)$, then
\[
 \code{incRev}(s\mathbin{+\!+}p)=t\mathbin{+\!+}p.
\]
\end{lemma}

\begin{proof}
For stripping, induct on $s$, generalizing the number of remaining matching
bits. At zero remaining bits the list is unchanged. Otherwise a matching head
decreases the counter and a nonmatching head does not; successful evaluation
ensures recursion never reaches an empty list with a positive counter.
For increment, a false head ends the carry and a true head propagates it.
Successful evaluation excludes propagation past the end. Induction gives both
equations. Reversal yields the corresponding most-significant-first suffix
statements.
\end{proof}

Let \code{controlSuffix} run the stripping check for continue, and stripping
followed by checked increment for break. It rejects the other edge kinds.
The checked generic theorem is
\[
\begin{split}
 \code{controlSuffix}(e,S)=\Some(E)
 \quad\Longrightarrow\quad
 \forall P,R,\;&\code{execute}(e,P\mathbin{+\!+}S,R)\\
 &=\Some(P\mathbin{+\!+}E,R).
\end{split}
\]
It is not sufficient merely to observe equality for one concrete caller
prefix. In a break instruction, both the stripping and the subsequent carry
must stay inside the chosen suffix.

\subsection{The combined quantified theorem}

A descriptor records a runtime-template index, a numerical entry state $q$,
a nonempty suffix $S$, a result suffix $E$, and its source edge $e$. A finite
kernel check verifies all of the following for each of the 35 descriptors:

\begin{itemize}
\item the descriptor selects the intended literal runtime template;
\item its entry is the suffix port associated with $q$ and $S$;
\item its control action passes \code{controlSuffix};
\item $|E|\leq|S|$, and its right endpoint has exactly the required normalized
suffix, gap, marker, and unchanged register box;
\item the descriptors enumerate all template indices $0,\ldots,34$.
\end{itemize}

Write $\mathcal R(v)=\code{sepsR}(v)$, with concatenation of Boolean words
written by juxtaposition. For each admitted descriptor, define
\[
\begin{aligned}
 A(q,P,S,n,v)
   &=z_c\bigl(q,\rev(\operatorname{init}S)\rev(P)1,
                   \operatorname{last}(S)0^{n+16}1\mathcal R(v)\bigr),\\
 B(P,S,E,n,v)
   &=z_c\bigl(84,1,PE0^{n+16+|S|-|E|}1\mathcal R(v)\bigr).
\end{aligned}
\]

\begin{theorem}[All 35 complete control contracts]
\label{thm:contracts}
Assume the existing \code{BoxLemmas}, Home rows and Home boundary contracts
for the maintained \code{tm215}. For every one of the 35 descriptors, every
prefix $P$ accepted by the existing scanner, every $n\in\mathbb N$, every
register list $v$ with $|v|\geq10$, and every source register function $R$,
\[
 \ReachOne_{\code{tm215}}\bigl(A(q,P,S,n,v),B(P,S,E,n,v)\bigr)
\]
and
\[
 \code{execute}(e,PS,R)=\Some(PE,R).
\]
\end{theorem}

\begin{proof}
Extract the literal guarded replay from the checked template list and apply
Corollary~\ref{cor:home}. The alignment check identifies the initial symbolic
configuration. Its right-normal-form equality, transported through the
existing \code{flatT_norm}, identifies the final \emph{whole} right tape for
every $n,v$. The checked length inequality justifies natural-number gap
arithmetic. Apply the source-action lemma to the checked \code{controlSuffix}
equation. The descriptor-index check ensures that this supplies every template,
not merely a selected example. The actual checked theorem is
\code{KFramedControl215.control_contract}.
\end{proof}

Only the right endpoint is normalized. One must not use \code{normL} to trim
trailing blanks before appending the frame. For example, $[0]$ and $[]$ can
represent the same blank unframed exterior, while appending $[1]$ gives
$[0,1]$ and $[1]$, which place the marker at different distances. The frame
proof deliberately preserves the complete symbolic left fragment.

As an example, the descriptor with entry state 13 and control instruction
\code{CC2} has $S=111$ and $E=1$. Its local body takes eight symbolic steps.
For every allowed prefix $P$, continue twice takes $P111$ to $P1$, leaves
registers unchanged, and increases the represented gap from $n+16$ to $n+18$.
Choosing a prefix ending in zero is permitted. The theorem is independent of
that example and quantifies over all admissible prefixes.

\subsection{Match the actual site gap}

The maintained site format uses gap $a+(64-|Q|)$, where the consumed source
word $Q$ has length at most 48. If $Q=PS$, instantiate the template parameter
with
\[
 n=a+(48-|Q|).
\]
Then $n+16=a+(64-|Q|)$. Since $|E|\leq|S|$, the output gap equals
\[
 n+16+|S|-|E|=a+(64-|PE|).
\]
These are equalities in natural-number arithmetic; the length bounds are
essential when simplifying subtraction. The template's source action also
has no prefix carry or marker overflow. The independent occurrence audit
checks the exact action, route, decomposition, scanner condition, and gap
delta at all 1,422 assigned sites.

\section{Use semantic coverage as the consumer interface}
\label{sec:integration}

The maintained \code{KSim.Covered} currently asks that each decoded leaf satisfy
the bounded Boolean \code{checkEdge}. The generic endpoint consumes that exact
predicate \cite{ksim,endpoint}. Theorem~\ref{thm:contracts} is a semantic
reachability theorem. It does not imply that the old executor, with its fixed
fuel, returns \code{true}. A completeness theorem of that kind is neither
available nor generally valid for arbitrary bounded execution.

Consequently, the new theorem cannot simply be substituted for
\code{KSites215.covered}. The correct interface is an edge contract:
\[
 \operatorname{EdgeContract}(M,Q,e),
\]
which requires the same path/register bounds and the appropriate positive
whole-tape executions, with separate zero and nonzero decrement branches and
a separate halt clause. The already developed \code{semantic-coverage}
proposal supplies a semantic consumer, a sound adapter from \code{checkEdge},
and mixed coverage \cite{semantic}. Reuse and rebase that work; do not
reimplement the simulation induction independently for each template family.

The intended proof has two producers of the same contract:
\begin{enumerate}
\item \emph{Literal producer.} The existing checked executor proves the
remaining 4,754 outcomes: 1,900 increments, 2,850 decrement branches, three
control fallbacks, and one halt.
\item \emph{Framed producer.} The decoder-prefix theorem reaches the exact
primitive port; Theorem~\ref{thm:contracts} performs the local control action
and positive Home return. Exact finite binding data instantiate it at the
other 1,422 outcomes.
\end{enumerate}

The consumer proves total source simulation from these contracts. A forward
step must have positive machine duration, so a halting machine cannot
simulate infinitely many source steps in zero machine time. The unchanged
Home theorem supplies this positivity for the new control family. The halt
clause continues to identify an actual terminal state, rather than an
exhausted symbolic evaluator.

\subsection{Precise remaining assembly obligations}

The following are still required before claiming that the shorter proof has
replaced the production endpoint:

\begin{enumerate}
\item Prove in the kernel that each selected source occurrence has its exact
consumed path, source action, and decoder route. Numerical equality of an
entry state is not itself a \code{RoutedEntry} proof: different decoder roles
can share a physical state.
\item Assemble the 1,422 template instances, using the positive concrete
decoder prefix and the exact input/output shape and gap rewrites. The new
generic source-action theorem removes the need to reprove prefix preservation
at every instance.
\item Retain the three exceptional control cases and all other outcomes through
the existing sound checker. Establish total semantic coverage of all 4,751
leaves and 6,176 outcomes, including initial entry and the literal halt leaf.
\item Add or rebase a semantic endpoint adapter while preserving the original
endpoint type and compatibility clients. Old Boolean coverage should remain
a sufficient way to use the generic endpoint.
\item Rebuild affected dependencies and the exact endpoint clients using the
production versions, inspect assumption reports, and separately measure
source size, generated data, object size, time, and memory.
\end{enumerate}

The new package does not include a completed full-site formal assembly or a
fresh ZFC endpoint build. Its independent Python audit supplies exact finite
data validation, not a substitute for these kernel assembly obligations.
The earlier round13 prose mentions an all-leaf residual packet under another
proposal path; this review does not treat that historical claim as a fresh
binding of sources that are absent at the review pin.

\section{A second component: extensional source rewriting in Rocq}
\label{sec:extensional}

\subsection{The exact local change}

The earlier CFG investigation found that the source programs differ in the
order of increment-only pairs at two shared \code{dagPerm} nodes
\cite{cfgproposal}:
\[
 \begin{array}{c|c|c}
 \text{node}&\text{original body}&\text{commuted body}\\\hline
 188&(\operatorname{Inc}6,\operatorname{Inc}7)&(\operatorname{Inc}7,\operatorname{Inc}6)\\
 190&(\operatorname{Inc}6,\operatorname{Inc}1)&(\operatorname{Inc}1,\operatorname{Inc}6).
 \end{array}
\]
The proposed program \code{dagPermComm} changes exactly these two node
definitions; its root and node count are unchanged. It is an intermediate
source presentation, not a new binary machine.

Two source templates account for 256 expanded swaps in the prior structural
analysis. This is precisely where shared-DAG reasoning helps: the semantic
proof should depend on the two changed node bodies, not on every expanded
occurrence. The existing \code{LocalProgram} proposal has a contextual
transfer principle with literal register-function equality \cite{localprogram}.
That equality is stronger than needed for these reordered updates.

\subsection{Use pointwise register equality}

A register file in the Rocq semantics is a function $R:\mathbb N\to\mathbb N$
\cite{rm}.
Define
\[
 R\RegEq S\quad\Longleftrightarrow\quad\forall i,\;R(i)=S(i).
\]
Pointwise equality lets us prove update commutation without converting this
statement into equality of function objects. That conversion would ordinarily
invoke functional extensionality in this proof style. The delivered source
transport never needs it.

\begin{lemma}[Increment commutation]
Let $I_a$ increment register $a$. For all $a,b,R$,
\[
 I_b(I_a(R))\RegEq I_a(I_b(R)).
\]
The lemma even permits $a=b$.
\end{lemma}

\begin{proof}
If $a=b$, both sides are the same nested increment. Otherwise fix an index
$i$. If $i=a$, both sides have value $R(a)+1$; if $i=b$, both have value
$R(b)+1$; otherwise both have $R(i)$. This is an indexwise argument using only
decidable equality of natural numbers. No function-equality axiom occurs.
\end{proof}

The actual structured semantics uses judgments
\[
 \operatorname{run}_p(e,R,x,R'),
\]
where $x$ is the normal, break, continue, or halt exit (with the existing
depth parameters where applicable). The new one-direction node contract is
\[
\begin{split}
 \operatorname{NodeTransport}(p,q,n)\;:\Longleftrightarrow\;
 \forall R,x,R',\;&\operatorname{run}_p(\operatorname{Node}n,R,x,R')\\
 \Longrightarrow\forall S,\;&R\RegEq S\Longrightarrow
 \exists S',\operatorname{run}_q(\operatorname{Node}n,S,x,S')\wedge R'\RegEq S'.
\end{split}
\]
The structured exit is preserved exactly. Only the register representation is
related pointwise.

\begin{definition}
\code{ClosedExt p q} says that for every node $n$, either its two child edges
are literally unchanged, or \code{NodeTransport p q n} holds.
\end{definition}

\begin{theorem}[Extensional contextual transport]
\label{thm:transport}
If \code{ClosedExt p q}, then every finite structured run of $p$ transports
to a run of $q$ with the same exit and pointwise related output registers,
starting from any pointwise related input registers.
\end{theorem}

\begin{proof}
Induct on the derivation of the finite structured run. Primitive increments
preserve pointwise equality under update; decrement branches agree because
the tested register values agree. Primitive break, continue, and halt retain
the same registers and exit.

At a node, inspect the closure condition. For a changed node, use its supplied
whole-node transport contract. For an unchanged node, invoke the induction
hypotheses on the first child, second child, and any loop-restart derivation
required by the actual run constructor. Pass the intermediate pointwise
register relation to the next induction hypothesis. Rebuild the same node
constructor using equality of the child-edge pair.

This handles all structured exit constructors, including the continuation
case that restarts the node. There is no acyclicity or node-rank hypothesis:
the induction is on the finite execution evidence, not on the source graph.
\end{proof}

\begin{corollary}[Halting equivalence]
If \code{ClosedExt p q}, \code{ClosedExt q p}, and the root edges agree, then
for every initial register file $R$,
\[
 \operatorname{iter\_halts}(p,R)
 \quad\Longleftrightarrow\quad
 \operatorname{iter\_halts}(q,R).
\]
\end{corollary}

\begin{proof}
Transport each finite root run with Theorem~\ref{thm:transport}. For a normal
root exit, feed the pointwise output relation into the induction hypothesis
for the remaining iterations; for a halt exit, finish. Apply this argument in
both directions, starting with $R\RegEq R$.
\end{proof}

\subsection{Checked concrete result and its limitation}

Both changed nodes contain exactly two increments. Their structured runs must
exit normally with the corresponding pair of updates. The increment lemma
proves the node contract in either direction. All other nodes are literally
unchanged. The checked result is
\[
 \forall R,\quad
 \code{iter_halts dagPerm R}
 \quad\Longleftrightarrow\quad
 \code{iter_halts dagPermComm R}.
\]
The package also supplies transport of arbitrary structured edge runs in both
directions, preserving every exit and pointwise output registers.

The generic and concrete components were freshly compiled with Coq 8.18.0,
using the actual pinned \code{RM.v} and current \code{DagPerm.v}. All six
printed assumption reports are closed under the global context. Recursive
\code{coqchk} of the concrete result passes. This is a focused compatibility
check, not a claim that the full production Rocq 9.1.0 endpoint has been
recompiled with Coq 8.18.0.

\textbf{What this does not yet prove.} The previous Python experiment reports
ordered-CFG equality of \code{p215} with the two-node-commuted source. The new
Rocq theorem does not turn that experiment into a kernel theorem. In
particular, it does not eliminate the 1,426-entry normalized source
correspondence certificate by itself.

To complete that route, certify flattening or a canonical CFG checker that
preserves initial entry, every decrement branch, explicit halt, and all
eliminated control transfers. Pure-control cycles must be rejected or
represented as divergence. Increment-only cycles must not be replaced by a
finite increment vector. A generic certificate about finite elimination
traces or a decreasing rank can establish that normalization does not hide
infinite zero-progress execution. The earlier analysis already identifies
these obligations; the present result closes the local source-commutation
part while leaving flattening explicit.

\section{Validation and artifact interpretation}
\label{sec:validation}

The archive separates three kinds of evidence:
\begin{description}
\item[Fresh kernel checks.] The Lean frame rule, guarded replay, Home composition,
source-action framing, all 35 literal runtime checks, descriptor alignment,
and combined universal control contracts. The Rocq pointwise contextual
transport and exact two-node source rewrite, with recursive checking;
also the separately checked Rocq local frame rule and 35 local replays,
with the precise recursive-check boundary stated above.
\item[Fresh finite audits.] An independently authored Python checker binds the
35 templates and 1,422 occurrences to the pinned machine, maintained decoder,
and complete source expansion. It checks all 444 local steps, full state-label
translation, source-action guards, right-tape arithmetic, and the exact three
fallbacks. It contains explicit negative examples at the unsafe boundary.
\item[Historical context.] The repository's endpoint/checkpoint receipts and
earlier proposal findings. These explain what was already accomplished and
where the new components fit. They are not relabelled as this delivery's
fresh endpoint checks.
\end{description}

The source tree contains the unchanged focused dependencies required for
offline component replay. Provenance manifests distinguish those inherited
files from new modules and bind each to exact bytes. Compiler binaries and
compiled proof objects are not shipped as substitutes for source checking.
Raw successful logs, exact compiler versions, commands, hashes and results
are retained. Failed development attempts, where retained, are separate from
the accepted fresh replay; a printed \code{sorryAx} in a failed attempt is not
an accepted theorem report.

The principal new Lean declarations are:
\begin{itemize}
\item \code{KLeftFrame.sstep_natural}, \code{runLocal_sound}, and \code{local_home};
\item \code{KControlFrame.controlSuffix_frame};
\item \code{KLeftFrame215.all_checked} and \code{template_home};
\item \code{KFramedControl215.all_aligned}, \code{descriptor_indices}, and
\code{control_contract}.
\end{itemize}
For Rocq, \code{run_transfer_ext} and \code{iter_iff_ext} supply generic
transport. The declarations \code{swap_increment_node} and
\code{dagPerm_comm_iter_iff} supply the concrete application.
The runtime component adds the guarded frame
and local replay declarations described in Section~\ref{sec:frame}.

See the archive's \code{README.md} for exact replay commands and the final
machine-readable verification summary for the accepted compilation scope.
The integrity verifier checks archive bytes; the proof compilers and recursive
checker establish the corresponding formal claims. The inventory and finite
audit scripts do not claim to establish an internal ZFC derivation.

\section{Integration plan and meaningful measurements}

\begin{longtable}{@{}p{0.12\textwidth}p{0.36\textwidth}p{0.44\textwidth}@{}}
\toprule
Stage & Change & Completion criterion\\
\midrule\endhead
1 & Import the new frame and control-action lemmas & Fresh focused checks on
the pinned production toolchains; existing executor definitions unchanged\\
2 & Rebase semantic coverage with compatibility adapters & Old literal checker
still produces the same contracts; old endpoint API clients still typecheck\\
3 & Assemble 1,422 exact template occurrences & Kernel-bound decoder roles,
paths, source actions, gap equations, and prefix premises; no missing outcomes\\
4 & Compose mixed semantic coverage & All 4,751 leaves and 6,176 outcomes
covered, with three control fallbacks and a separate terminal case\\
5 & Rebuild the exact K215 endpoint & Same literal-table hashes, same final
theorem types, allowed assumption reports, no incomplete providers\\
6 & Develop certified source flattening & Compose with the checked two-node
Rocq equivalence; preserve initial entry, branch order, halt and divergence\\
\bottomrule
\end{longtable}

For stages 3--5, benchmark the baseline and candidate on the same toolchain,
machine, resource limits, and fresh or equivalently warmed object conditions.
Report at least:
\begin{enumerate}
\item maintained handwritten source and generated source separately;
\item serialized finite data and compiled object sizes;
\item elaboration/kernel-check time and peak memory for the runtime family;
\item affected endpoint build time, with dependency reuse explicitly identified;
\item number of independent expensive executions and number of total bindings.
\end{enumerate}

The runtime contract module adds generic proof infrastructure. Its first
adoption may increase source lines even while reducing repeated computation.
Do not report the 22.46\% obligation reduction as a 22.46\% source reduction
or time reduction. Similarly, a small final wrapper does not establish that
the underlying checked dependency closure is smaller. The delivered components
make a concrete comparison possible; they do not predetermine its result.

For Rocq, pointwise transport is useful independently of the two-node
normalizer. It permits future register-permuting or update-reordering local
contracts without importing functional extensionality. The exact current
closed endpoint should remain the reference assumption profile.

\section{Further research questions}

\begin{enumerate}
\item \textbf{Can the canonical CFG be computed and justified from the compact
source DAG?} The two-node commutation theorem now supplies the local rewrite.
A certified extraction theorem could remove separately maintained path-pair
correspondence data. It must explicitly handle control-only and increment-only
cycles, malformed walks, initial entry, and the true halt leaf.
\item \textbf{Can the same frame principle cover increments and decrements?}
They account for 4,750 guarded outcomes, substantially more than the control
family. The likely abstraction combines unchanged prefix with the existing
register effect and a parametric gap update. Decrement's two branch premises
and the extra spare block produced by increment must remain explicit.
\item \textbf{Can a small checker discover the minimal protected suffix?}
A reflected footprint analysis could find the first boundary that a local
trace must not cross. Its correctness should imply the guarded frame rule,
and its rejection path should fall back to the original checker. This would
reduce hand-maintained occurrence assignments without trusting the discovery
algorithm.
\item \textbf{What is the best compiled representation of the 35 descriptors
and 1,422 bindings?} A compact indexed array, source-DAG annotation, or shared
proof tree may reduce evaluation and source maintenance differently. Compare
actual object size and checked execution cost; printed line count alone is a
poor guide.
\item \textbf{How much of the symbolic-progress flag can remain local?}
The present proof needs no positivity from the guarded fragment because Home
supplies it. More generally, require progress at a designated observable
boundary instead of rebuilding positive duration bookkeeping for each
administrative symbolic operation.
\item \textbf{Can extensional contextual transport be shared across both proof
lanes as a specification?} The languages should retain independent small
proofs, but agree on the exact exit relation, pointwise register relation,
node replacements, and certificate data. This avoids a new trusted
cross-language translator.
\item \textbf{What is required for the separate internal ZFC endpoint?}
The existing proof-assistant endpoint and the reviewed ordinary mathematical
argument are distinct from an encoded derivation in ZFC syntax. Any future
internalization should expose its proof-code translation and arithmetic
formalization instead of treating the external theorem as that object.
\end{enumerate}

\section*{Attribution}
\addcontentsline{toc}{section}{Attribution}

The 35-template discovery and empty-left decomposition come from the existing
Round 13 left-frame proposal. The two-node source comparison comes from
the third K215 simplification proposal. The contextual source-transfer design builds on
the \code{common-macros} proposal and the existing Wade register-machine
semantics. The symbolic executor, zipper model, Home routine, source
calibration, K215 construction, and original Lean/Rocq proofs are inherited
from the named contributors and upstream projects documented in ProveIt.
New work here comprises the checked frame completion, source-action framing,
35 combined Lean contracts, independently checked Rocq local frames and
replays, pointwise Rocq transport, concrete commutation proof,
independent audits, and this integration analysis. Original notices and
licenses accompany the reused source files.

\begin{thebibliography}{20}
\small
\bibitem{verification} ProveIt contributors.
\href{\baseurl Computability/BusyBeaver/Research/ZFC/k225-formal/verification/k215/README.md}{Formal correctness of the exact 215-state machine}.
Review snapshot specified in Section~\ref{sec:baseline}.

\bibitem{candidate} ProveIt contributors.
\href{\baseurl Computability/BusyBeaver/Research/ZFC/candidate-k215/README.md}{K215: 215 working states for ZFC consistency}.

\bibitem{simplifications} ProveIt contributors.
\href{\baseurl Computability/BusyBeaver/Research/ZFC/k225-formal/verification/k215/simplification/README.md}{Shorter checked K215 proofs}.

\bibitem{phase15} ProveIt contributors.
\href{\baseurl Computability/BusyBeaver/Research/ZFC/k225-formal/verification/k215/simplification/phase15/README.md}{Fifteenth K215 proof-simplification checkpoint}.

\bibitem{leftproposal} ProveIt proposal archive.
\href{\baseurl Computability/BusyBeaver/Research/ZFC/k225-formal/verification/k215/simplification/proposals/round13-left-frame/runtime/README.md}{Round 13 runtime: factor at the first empty-left Home boundary}.
Includes the original uncompiled contract sketch and template data.

\bibitem{cfgproposal} ProveIt proposal archive.
\href{\baseurl Computability/BusyBeaver/Research/ZFC/k225-formal/verification/k215/simplification/proposals/k215-simplification-3/rocq/CFG_REPORT.md}{Two local commutations instead of a general source certificate}.

\bibitem{localprogram} ProveIt proposal archive.
\href{\baseurl Computability/BusyBeaver/Research/ZFC/k225-formal/verification/k215/simplification/proposals/common-macros/rocq/src/KLocal/LocalProgram.v}{LocalProgram: whole-node source transfer}.

\bibitem{semantic} ProveIt proposal archive.
\href{\baseurl Computability/BusyBeaver/Research/ZFC/k225-formal/verification/k215/simplification/proposals/semantic-coverage/README.md}{K215: semantic coverage and a shared implication core}.

\bibitem{ksim} ProveIt contributors.
\href{\baseurl Computability/BusyBeaver/Lean/BusyBeaver/ZFC/KSim.lean}{Maintained KSim: decoder sites and register-machine simulation}.

\bibitem{endpoint} ProveIt contributors.
\href{\baseurl Computability/BusyBeaver/Lean/BusyBeaver/ZFC/KEndpoint.lean}{Maintained KEndpoint: generic endpoint composition}.

\bibitem{runtime} ProveIt contributors, adapting the K230 and m254 runtime proofs.
\href{\baseurl Computability/BusyBeaver/Lean/BusyBeaver/ZFC/KRuntimeCore.lean}{KRuntimeCore: symbolic execution and soundness}.

\bibitem{home} ProveIt contributors.
\href{\baseurl Computability/BusyBeaver/Lean/BusyBeaver/ZFC/KHome.lean}{KHome: shared all-gap Home scan}.

\bibitem{rm} K230 and upstream contributors.
\href{\baseurl Computability/BusyBeaver/Research/ZFC/k225-formal/verification/k215/simplification/proposals/zfc215-simplification/source_proof/deps/K230/RM.v}{Pinned Wade register-machine semantics, RM.v}.
The archive retains the exact inherited source and attribution.
\end{thebibliography}
\end{document}
